% paper-en.tex —— paper.md 的排版渲染 (英文), 供预览与投稿前校对。
%
% paper.md + paper.bib 是 JOSE 投稿的**唯一源**; 本文件只负责把同一份内容
% 排成 PDF。改动正文时请改 paper.md, 再同步这里 —— 两处内容必须一致。
%
% 编译: xelatex paper-en && bibtex paper-en && xelatex paper-en && xelatex paper-en
\documentclass[11pt,a4paper]{article}

\usepackage{fontspec}
\usepackage[margin=2.4cm]{geometry}
\usepackage[authoryear,round]{natbib}
\usepackage{titlesec}
\usepackage{enumitem}
\usepackage{xcolor}
\usepackage[hidelinks]{hyperref}

\setmainfont{Times New Roman}
\setlist{nosep,leftmargin=1.4em}
\titlespacing*{\section}{0pt}{1.1em}{0.45em}
\titlespacing*{\subsection}{0pt}{0.8em}{0.35em}
\titleformat{\section}{\normalfont\large\bfseries}{\thesection.}{0.5em}{}

\newcommand{\doilink}[1]{\href{https://doi.org/#1}{#1}}

\title{\vspace{-1.6cm}\bfseries From Bits to Bootstrapping:\\[3pt]
A Bilingual, Hands-On Toolkit for Teaching nsF5 Steganography\\[3pt]
and Its Detection}
\author{Fengjie Lin\thanks{Qilu University of Technology.
  ORCID: \texttt{0000-0000-0000-0000} (placeholder --- to be filled before submission).}}
\date{13 September 2026}

\begin{document}
\maketitle
\vspace{-1.2em}

\section*{Summary}
\addcontentsline{toc}{section}{Summary}

\texttt{nsf5stego} is an open-source teaching and research toolkit for image
steganography and steganalysis. It provides a complete, readable implementation
of the nsF5 algorithm --- Hamming syndrome matrix coding combined with wet paper
coding \citep{westfeld2001f5,fridrich2005wet} --- together with SHA-256
content-keyed embedding, blind steganalysis by chi-square and RS analysis
\citep{westfeld1999,fridrich2001rs}, and supervised detection with
gradient-boosted classifiers \citep{ke2017lightgbm}. Around that core it ships a
bilingual handbook (Chinese and English, 66 and 72 pages), ten runnable
notebooks, an interactive browser laboratory, and a desktop GUI.

The design premise is that the three ideas which actually make modern
steganography work --- syndrome coding, the wet paper model, and the statistical
footprint that embedding leaves behind --- are learned by building them and then
attacking what you built, not by reading about them. Every algorithm in the
toolkit is therefore exposed as an object the learner can inspect and modify at
intermediate steps, rather than as a library call that returns a stego image.

\section*{Statement of need}

Steganography teaching material tends to sit at one of two extremes. Textbooks
treat it as mathematics: they derive the syndrome equation and the wet paper
model rigorously, but leave the reader without a running pipeline
\citep{fridrich2009book,provos2003hide}. Course modules that do provide code tend
to stop at LSB substitution, which is the one technique that the field moved past
two decades ago precisely because it is trivially detectable. The result is that
a student can often explain the Hamming bound while being unable to say what wet
paper coding buys, why F5 shrinks the message, or what an embedding efficiency of
4 bits per change actually costs.

There is a second, subtler gap. Detection performance is overwhelmingly reported
on curated benchmark corpora built for the purpose \citep{bas2011boss}, where
source images are partitioned and cover/stego pairs are carefully constructed
\citep{pevny2010spam,fridrich2012rich}. Learners who train on their own
photographs, or who split cover/stego pairs from the same source image across
training and test sets, obtain numbers that look excellent and mean nothing.
\texttt{nsf5stego} is built so that this mistake is difficult to make and easy to
see: the dataset builder emits a \texttt{photo\_id} for every sample, and every
evaluation path groups by it by default. An audit of the project's own
experiments found a published comparison in which a 53-dimensional result had in
fact been computed from 11 features --- the exact class of error the toolkit now
guards against with hard assertions rather than silent fallbacks.

\section*{Target audience and instructional context}

The material targets advanced undergraduates and beginning graduate students in
security, signal processing, or applied machine learning. It assumes introductory
Python and no prior exposure to information hiding; the handbook opens with pixels
and binary representation before introducing anything steganographic. It has been
designed as a one-semester elective or as the implementation backbone of a
capstone project, with a twelve-week quick-start path and a six-to-twelve-month
deep track for students continuing into research.

\section*{What the toolkit contains}

The package is organised so that theory, implementation, and verification stay
adjacent:

\begin{itemize}
\item \textbf{A reference implementation} of nsF5 with binary Hamming codes
  $[n = 2^p - 1,\, k,\, 3]$, magnitude-decrement modification, and a wet paper
  solver that resolves dry positions through progressively deeper GF(2) searches.
  Optional C++ accelerators are compiled from source on Windows, macOS, and
  Linux; when absent, equivalent pure-Python paths run automatically, so
  notebooks and containers work with no build step.
\item \textbf{Attacks on what you just built.} Chi-square and RS analysis run
  live against the learner's own embedded image, and the browser laboratory draws
  the mask of changed pixels next to the detector verdict, so the connection
  between ``these LSBs moved'' and ``this detector fired'' is visible rather than
  asserted.
\item \textbf{An interactive browser laboratory} covering bit-plane extraction
  and weighted restacking, an embed/decode round trip with live detection, a
  walkthrough of syndrome coding, a wet-paper dry-position solver, a side-by-side
  comparison of naive LSB against the project's real nsF5, an ML decision-threshold
  playground, and a payload scan computed from the project's own statistics.
\item \textbf{Ten notebooks} that execute as Colab or Jupyter sessions, each
  mapped to a handbook chapter \citep{perez2007ipython}.
\item \textbf{A narrated video series} in Chinese (eleven chapters, about 36
  minutes total) for students who encounter the written material before the
  mathematics it depends on.
\item \textbf{A dataset builder} that produces cover/stego corpora with
  per-source-image identifiers, and a downloader for the BOSSbase 1.01 benchmark
  \citep{bas2011boss}.
\end{itemize}

\section*{Learning outcomes}

By working through the material, a student should be able to: implement and
invert Hamming syndrome matrix coding, and compute its embedding efficiency;
explain why F5 shrinks the message and how nsF5 removes that shrinkage at the
cost of a denser search; derive and solve the wet paper system for a given set of
dry positions; implement chi-square and RS detectors and interpret their
disagreement; build a supervised steganalysis classifier with grouped
cross-validation and report it with confidence intervals; and --- the outcome the
project most cares about --- recognise when a reported detection result is an
artefact of evaluation design rather than a property of the detector.

This last point is taught through a worked negative example. A CNN
reimplementation in the repository attains a chance-level AUC on both validation
\emph{and training} data, which the toolkit's reporting makes immediately legible
as non-convergence rather than as evidence about convolutional networks, and
which would otherwise have been publishable as a misleading comparison.

\section*{Relationship to existing resources}

General introductions \citep{provos2003hide} and the standard monograph
\citep{fridrich2009book} supply the theory but not a pipeline. Attack
implementations exist primarily as research code accompanying papers
\citep{pevny2010spam,fridrich2012rich}, which is optimised for throughput rather
than for inspection. Deep-learning steganalysis is well served by reference
architectures \citep{xu2016xunet,ye2017yenet}, but those presuppose the classical
material this toolkit supplies. \texttt{nsf5stego} differs by covering the whole
arc --- pixel representation, embedding, keying, detection, classifier evaluation
--- in one inspectable codebase, and by treating the evaluation protocol itself as
part of the curriculum.

\section*{Verification and current use}

The software is Apache-2.0 licensed and archived with a DOI. Twenty-four automated
tests run under continuous integration on Ubuntu, macOS, and Windows, alongside a
container build and a browser regression suite that exercises the interactive
laboratory. The C++ and Python embedding paths are verified to interoperate, and
the project documents where they legitimately differ rather than asserting
equivalence.

The material has not yet been adopted in a formal course; the toolkit's
verification evidence at this stage concerns the software's correctness, not its
pedagogical effectiveness. The author welcomes collaboration with instructors
willing to trial it, and would particularly value reports of the misconceptions
students actually hold, which is the evidence the project currently lacks.

\section*{Acknowledgements}

The toolkit builds on the published steganography and steganalysis literature
cited above, and on the NumPy and Matplotlib scientific Python stack
\citep{harris2020numpy}. BOSSbase 1.01 is used as the reference benchmark
\citep{bas2011boss}.

\bibliographystyle{plainnat}
\bibliography{paper}

\end{document}
